\documentclass[11pt]{article}
\usepackage[margin=1in]{geometry}
\usepackage{amsmath,amssymb,amsthm}
\usepackage{graphicx}
\usepackage{hyperref}

\newtheorem{theorem}{Theorem}
\newtheorem{proposition}{Proposition}
\newtheorem{lemma}{Lemma}
\newtheorem{conjecture}{Endpoint Lemma}
\newtheorem{remark}{Remark}

\newcommand{\E}{\mathbb E}
\newcommand{\Pp}{\mathbb P}
\newcommand{\R}{\mathbb R}
\newcommand{\cB}{\mathcal B}
\newcommand{\norm}[1]{\left\lVert #1\right\rVert}
\newcommand{\Lq}{L^q(S)}
\newcommand{\sfq}{\mathsf{SF}_q}
\newcommand{\lesssimq}{\lesssim_q}
\newcommand{\simeqq}{\simeq_q}

\title{A Conditional Reduction for the \(L^q\), \(p=1\) Endpoint}
\author{Agent \texttt{agent\_lane\_19}}
\date{2026-06-16}

\begin{document}
\maketitle

\section*{Status}

This is a conditional reduction packet for Remark 12.10 of Yaroslavtsev's
paper \emph{Local characteristics and tangency of vector-valued martingales}
\cite{Yaroslavtsev2019}.  The earlier packet
\[
\texttt{1907.11588\_Lq\_general\_local\_characteristic\_domination}
\]
proves the \(L^q(S)\) case for \(1<p,q<\infty\).  This packet isolates the
remaining endpoint \(p=1\), \(1<q<\infty\), to one explicit square-function
comparison.

The endpoint comparison is not proved here.  Thus the packet is conditional:
if the endpoint square-function aggregate comparison in
Lemma \ref{lem:endpoint-square} is true for a fixed \(q\), then the full
\(L^q(S)\)-valued characteristic-domination conclusion at \(p=1\) follows for
that \(q\).

\section*{Source Problem}

The source is arXiv:1907.11588, PDF page 105, Remark 12.10.  The source states
that the general local-martingale case is open and reduces the missing issue
to an independent-increment comparison.

\begin{center}
\includegraphics[width=\textwidth]{figures/open_problem_crop.png}
\end{center}

The continuation states the aggregate-domination formulation and notes that
even the symmetric case was not known to the author.

\begin{center}

\includegraphics[width=\textwidth]{figures/open_problem_continuation_crop.png}
\end{center}

\section*{Conditional Dependency}

For a finite independent family \(z=(z_i)\) of \(L^q(S)\)-valued random
variables, put
\[
  \sfq(z)
  =
  \E\left\|
      \left(\sum_i |z_i|^2\right)^{1/2}
    \right\|_{L^q(S)}.
\]
Let
\[
   \mu_z=\sum_i \mathcal L(z_i)|_{\Lq\setminus\{0\}}
\]
be the aggregate occurrence measure.

\begin{conjecture}[Endpoint square-function aggregate comparison]
\label{lem:endpoint-square}
Let \(1<q<\infty\).  There is a constant \(C_q\) such that for every measure
space \(S\) and every two finite independent families \(x=(x_i)\),
\(y=(y_j)\) of \(L^q(S)\)-valued random variables,
\[
  \mu_x\le \mu_y
  \quad\Longrightarrow\quad
  \sfq(x)\le C_q\,\sfq(y).
\]
\end{conjecture}

This is the only unproved assertion used in the conditional theorem below.
It is true in Hilbert spaces by the Laplace-transform argument in the companion
Hilbert all-\(p\) packet, and the broad \(c_0\) counterexample shows that no
analog can hold in arbitrary Banach spaces.  Finite-dimensional regrouping
tests for fixed \(q<\infty\) found no growing obstruction.

\section*{Idea of the Proof}

At \(p=1\), the Dirksen--Yaroslavtsev Rosenthal norm still gives the upper
bound for independent \(L^q\)-valued sums, but their printed two-sided theorem
uses \(L^p\)-duality and is stated for \(p>1\).  Yaroslavtsev's later BDG
paper supplies the missing endpoint geometry in a different form: for UMD
Banach function spaces, and hence for \(L^q(S)\), the \(L^1\) size of a
martingale is equivalent to the \(L^1\) size of its lattice square function.

Thus the \(p=1\) discrete bottleneck is exactly the question whether this
lattice square function is stable under aggregate domination of the jump laws.
If Lemma \ref{lem:endpoint-square} holds, the discrete comparison follows at
once from the \(p=1\) Banach-function-space BDG theorem.  The passage from
discrete independent increments to general local martingales is then the same
total-characteristic reduction used in the existing \(1<p<\infty\) packet:
continuous parts are already controlled at \(p=1\), and jump parts are compared
only after accessible and quasi-left-continuous jumps are folded into the
total spatial compensator.

\section*{Discrete Endpoint Consequence}

\begin{proposition}[Conditional discrete \(p=1\) comparison]
\label{prop:discrete}
Assume Lemma \ref{lem:endpoint-square} for a fixed \(1<q<\infty\).  Let
\((\xi_i)_{i=1}^n\) and \((\eta_j)_{j=1}^m\) be independent mean-zero
\(L^q(S)\)-valued random variables.  If
\[
   \sum_i \Pp(\xi_i\in B)
   \le
   \sum_j \Pp(\eta_j\in B)
   \qquad(B\in\cB(\Lq),\;0\notin B),
\]
then
\[
   \E\max_{1\le k\le n}
      \left\|\sum_{i=1}^k \xi_i\right\|_{L^q(S)}
   \le
   C_q
   \E\max_{1\le k\le m}
      \left\|\sum_{j=1}^k \eta_j\right\|_{L^q(S)}.
\]
In particular,
\[
   \E\left\|\sum_i \xi_i\right\|_{L^q(S)}
   \le
   C_q
   \E\left\|\sum_j \eta_j\right\|_{L^q(S)}.
\]
\end{proposition}

\begin{proof}
The aggregate hypothesis is exactly \(\mu_\xi\le\mu_\eta\), hence Lemma
\ref{lem:endpoint-square} gives
\[
  \E\left\|\left(\sum_i|\xi_i|^2\right)^{1/2}\right\|_{L^q}
  \le
  C_q
  \E\left\|\left(\sum_j|\eta_j|^2\right)^{1/2}\right\|_{L^q}.
\]
Yaroslavtsev's Banach-function-space BDG theorem
\cite[Theorem 4.32 and Proposition 4.33]{Yaroslavtsev2018} says that for every
\(L^q(S)\)-valued martingale difference sequence \(d=(d_i)\),
\[
  \E\sup_N\left\|\sum_{i=1}^N d_i\right\|_{L^q}
  \simeqq
  \E\left\|
      \left(\sum_i |d_i|^2\right)^{1/2}
    \right\|_{L^q}.
\]
Both \((\xi_i)\) and \((\eta_j)\), in any fixed order, are martingale
difference sequences.  Applying the BDG equivalence to the two families gives
the maximal comparison.  The terminal comparison follows by dropping the
maximum.
\end{proof}

\begin{remark}
Dirksen and Yaroslavtsev's independent Rosenthal theorem
\cite[Theorem 2.1]{DirksenYaroslavtsev2017} explicitly keeps the upper bound
at \(p=1\).  Proposition \ref{prop:discrete} explains the missing lower side:
at the endpoint the right object is not a raw \(D_{q,q}\) or \(D_{1,q}\) term,
but the random lattice square function controlled by
Yaroslavtsev's \(p=1\) BDG theorem.
\end{remark}

\section*{Conditional General-Local Consequence}

\begin{theorem}[Conditional \(L^q\), \(p=1\) characteristic domination]
\label{thm:conditional-main}
Fix \(1<q<\infty\) and assume Lemma \ref{lem:endpoint-square}.  Let \(N\) and
\(M\) be \(L^q(S)\)-valued local martingales.  If \(N\) is characteristically
dominated by \(M\) in the sense of \cite[Definition 12.9]{Yaroslavtsev2019},
then
\[
  \E\sup_{t\ge0}\norm{N_t}_{L^q(S)}
  \le
  C_q
  \E\sup_{t\ge0}\norm{M_t}_{L^q(S)}.
\]
\end{theorem}

\begin{proof}
We first treat zero initial values.  Since \(L^q(S)\) is UMD, the source's
canonical and decoupled tangent machinery applies.  The continuous part is
already controlled for all \(0<p<\infty\) by Yaroslavtsev's continuous
characteristic-domination estimate \cite[Proposition 12.6]{Yaroslavtsev2019}.

It remains to compare the purely discontinuous contribution.  Conditional on
the local characteristics, the decoupled tangent jump martingales have
independent increments with deterministic total jump characteristics.  After
standard truncation of time and jump size, their jumps are finite independent
families whose aggregate occurrence measures are the corresponding total
spatial compensators.  Characteristic domination gives
\[
  \nu^N(\R_+\times B)\le \nu^M(\R_+\times B),
  \qquad B\in\cB(\Lq),\ 0\notin B.
\]
Therefore Proposition \ref{prop:discrete}, applied conditionally and then
passed to the limit by monotone/localization approximation, gives the \(L^1\)
maximal comparison for the jump tangent martingales.

The existing \(1<p<\infty\) packet verified the important structural point
needed here: the comparison uses only the total spatial jump compensator
\(\nu(\R_+\times\cdot)\), not a separate domination of accessible and
quasi-left-continuous parts.  Time labels and the accessible/qlc tag can be
erased for the identity jump integrand by conditional averaging over the
spatial jump value and by Radon--Nikodym restriction of decompositions.  The
same erasure is built into Lemma \ref{lem:endpoint-square}, which is stated
directly in terms of aggregate occurrence measures.

Combining the continuous and discontinuous estimates with the \(p=1\)
canonical-decomposition equivalence from Yaroslavtsev's BDG paper
\cite[Section 4.3]{Yaroslavtsev2018} yields the zero-start conclusion.

For nonzero initial values, Definition 12.9 gives
\(|\langle N_0,x^*\rangle|\le |\langle M_0,x^*\rangle|\) for all \(x^*\), hence
\(\norm{N_0}_{L^q}\le\norm{M_0}_{L^q}\).  Apply the zero-start result to
\(N-N_0\) and \(M-M_0\), and use the triangle inequality.
\end{proof}

\section*{Evidence and Failed Routes}

\paragraph{Hilbert endpoint.}
For Hilbert spaces, Lemma \ref{lem:endpoint-square} reduces to a scalar
fractional-moment comparison for
\((\sum_i \norm{\xi_i}^2)^{1/2}\).  The companion Hilbert all-\(p\) packet
proves this by the Laplace identity
\[
  \E S^\alpha
  =
  c_\alpha\int_0^\infty(1-\E e^{-tS})t^{-1-\alpha}\,dt,
  \qquad 0<\alpha<1,
\]
and the estimate
\[
  1-\prod_i(1-u_i(t))\simeq \min\{1,\sum_i u_i(t)\}.
\]
This argument uses only the aggregate law of the squared jump sizes.

\paragraph{\(c_0\) boundary.}
The broad Banach-space statement is false in \(c_0\): deterministic coordinate
jumps have supremum norm \(1\), while regrouping the same aggregate law by
random coordinates creates an occupancy maximum of order
\(\log N/\log\log N\).  This is exactly the obstruction one would expect to
disappear at fixed finite \(q\), where occupancy moments are finite constants
depending on \(q\).

\paragraph{Endpoint Rosenthal lower bound.}
A tempting route is to extend the Dirksen--Yaroslavtsev \(s_{p,q}\)-norm
equivalence to \(p=1\) by replacing \(D_{p,q}\) with \(D_{1,q}\).  Their paper
does state the independent upper estimate at \(p=1\), but the printed lower
bound uses \(L^p\)-duality and does not reach the endpoint.  Yaroslavtsev's
\(p=1\) Banach-function-space BDG theorem supplies a square-function lower
object, but I did not find a complete proof that the endpoint Rosenthal
sum/intersection norm is always controlled by this square function.

\paragraph{Scalar disjointification literature.}
Johnson--Schechtman, Montgomery-Smith, and the later Jiao--Sukochev--Xie--Zanin
disjointification theorems give sharp scalar formulas for
\[
  \left\|\left\|\sum_k f_k e_k\right\|_{\ell_q}\right\|_{L_p}
\]
for all \(0<p,q<\infty\).  These results strongly resemble the missing
endpoint lemma.  They do not immediately prove it because here each jump is an
\(L^q(S)\)-valued random function, so the dependence among spatial coordinates
inside one jump is part of the problem.

\section*{Finite Checks}

The script \texttt{code/grouping\_square\_simulation.py} randomly generates a
finite aggregate law in \(\ell_q^d\), partitions the same atoms into two
different independent families, and Monte Carlo estimates the two endpoint
square functions.  In several hundred random regroupings for
\[
q=1.1,1.2,1.5,1.8,2,3,5,10,20
\]
the largest observed ratios were below \(1.10\).  These checks are not proof;
they only support the claim that the \(c_0\) occupancy mechanism does not
persist at fixed finite \(q\).

\section*{Novelty and Literature Check}

Cheap run indexes were searched for \texttt{1907.11588}, \texttt{p=1},
\texttt{endpoint}, \texttt{Lq}, \texttt{square function},
\texttt{aggregate domination}, and \texttt{characteristic domination}.  Existing
nearby packets cover \(c_0\), Hilbert all finite \(p\), \(\ell_q\) at \(p=q\),
discrete \(L^q\) for \(1<p<\infty\), and full \(L^q\) general local martingales
for \(1<p<\infty\).  No local packet proves the \(L^q\), \(p=1\) endpoint.

Literature inspected during this push included Yaroslavtsev
\cite{Yaroslavtsev2019}, Dirksen--Yaroslavtsev
\cite{DirksenYaroslavtsev2017}, Yaroslavtsev's BDG paper
\cite{Yaroslavtsev2018}, Montgomery-Smith--Pruss
\cite{MontgomerySmithPruss1998}, and the scalar disjointification theorem of
Jiao--Sukochev--Xie--Zanin \cite{JiaoSukochevXieZanin2019}.  I found strong
nearby tools but no statement that directly proves Lemma
\ref{lem:endpoint-square}.

\section*{Human Review Recommendation}

Review the packet as a reduction, not as a solved endpoint theorem.  The
highest-value next task is to prove or disprove Lemma \ref{lem:endpoint-square}.
Promising proof routes are:
\begin{itemize}
\item a vector-valued Johnson--Schechtman disjointification theorem for the
mixed norm \(L^1(\Omega;L^q(S;\ell^2))\);
\item a Bernoulli--Poisson comparison for the \(L^q\)-lattice square function;
\item an endpoint Davis decomposition proving the \(p=1\) Rosenthal
sum/intersection norm is controlled by Yaroslavtsev's square function.
\end{itemize}

\begin{thebibliography}{9}

\bibitem{Yaroslavtsev2019}
I. S. Yaroslavtsev,
\emph{Local characteristics and tangency of vector-valued martingales},
arXiv:1907.11588, 2019.

\bibitem{DirksenYaroslavtsev2017}
S. Dirksen and I. S. Yaroslavtsev,
\emph{\(L^q\)-valued Burkholder--Rosenthal inequalities and sharp estimates for
stochastic integrals},
arXiv:1707.00109, 2017.

\bibitem{Yaroslavtsev2018}
I. S. Yaroslavtsev,
\emph{Burkholder--Davis--Gundy inequalities in UMD Banach spaces},
arXiv:1807.05573, 2018.

\bibitem{MontgomerySmithPruss1998}
S. J. Montgomery-Smith and A. R. Pruss,
\emph{A comparison inequality for sums of independent random variables},
arXiv:math/9811124, 1998.

\bibitem{JiaoSukochevXieZanin2019}
Y. Jiao, F. Sukochev, G. Xie, and D. Zanin,
\emph{Rosenthal's inequalities: \(\Delta\)-norms and quasi-Banach symmetric
sequence spaces},
arXiv:1910.12482, 2019.

\end{thebibliography}

\end{document}
